\documentclass[11pt,letterpaper]{article}

% Essential Packages
\usepackage[margin=1in]{geometry}
\usepackage{amsmath,amssymb}
\usepackage{graphicx}
\usepackage{circuitikz}
\usepackage{tikz}
\usetikzlibrary{shapes,arrows,positioning,calc}
\usepackage{booktabs}
\usepackage{multirow}
\usepackage{xcolor}
\usepackage{fancyhdr}
\usepackage{tcolorbox}
\usepackage{enumitem}
\usepackage{float}
\usepackage{listings}
\usepackage{hyperref}

% Colors
\definecolor{nlcblue}{RGB}{0,51,102}
\definecolor{alertred}{RGB}{204,0,0}
\definecolor{safgreen}{RGB}{0,128,0}

% Header/Footer
\pagestyle{fancy}
\fancyhead[L]{\textbf{NLC STEM Club UAV Project}}
\fancyhead[R]{\textbf{First Meeting Agenda}}
\fancyfoot[C]{\thepage}

% Title
\title{\Huge\textbf{First Meeting Agenda}\\
\Large Electrical Team Kickoff\\
\large NLC STEM Club - Electrical Systems Team}
\author{February 2026}
\date{}

\begin{document}

\maketitle
\thispagestyle{empty}
\newpage
\section{FIRST ELECTRICAL TEAM MEETING}

\begin{itemize}
  \item Agenda \& Leadership Guide
\end{itemize}

\begin{itemize}
  \item Fixed-Wing Environmental Monitoring UAV
  \item NLC STEM Club - Electrical Systems Team
  \item Meeting Duration: 2.5-3 hours
\end{itemize}

\section{THIS IS A PORTFOLIO-LEVEL PROJECT}

\begin{itemize}
  \item You are not building a hobby drone. You are executing
  \item aerospace systems engineering work that will demonstrate
  \item your competency to future employers and graduate programs.
\end{itemize}

\begin{itemize}
  \item Act like a systems engineer. Use professional terminology.
  \item Set the tone for excellence from day one.
\end{itemize}

\begin{itemize}
  \item PRE-MEETING PREPARATION (Your Checklist)
\end{itemize}

\begin{itemize}
  \item Print Documents (1 copy per attendee + 2 extras):
  \item This Meeting Agenda (19 pages)
  \item Power Budget Analysis (10 pages)
  \item Electrical System Block Diagram (7 pages)
  \item Detailed Wiring Diagram (16 pages)
  \item Safety Protocol (11 pages) - SIGNATURES REQUIRED
  \item Subsystem Assignment Sheet (see Section 9)
\end{itemize}

\begin{itemize}
  \item Gather Physical Materials:
  \item Whiteboard markers (3+ colors for block diagram)
  \item Component inventory (all electrical parts from parts list)
  \item Multimeter (for demonstrations)
  \item OVONIC 4S LiPo battery (fully charged, in safety bag)
  \item LiPo charger (B6 charger)
  \item Fire extinguisher (verify location, show to team)
  \item First aid kit (verify location)
  \item Pens for signatures
\end{itemize}

\begin{itemize}
  \item Prepare Workspace:
  \item Reserve room for 2.5-3 hours
  \item Set up tables for component inventory
  \item Label 4 stations: Power / Flight Electronics / Research / Integration
  \item Ensure electrical outlets available (for battery charging demo)
\end{itemize}

\begin{itemize}
  \item Mental Preparation:
  \item Review all printed documents thoroughly
  \item Rehearse opening statement (see Section 2)
  \item Prepare to answer questions confidently
  \item Remember: You are the expert. They are looking to YOU for leadership.
\end{itemize}

\section{MEETING STRUCTURE OVERVIEW}

\section{TIME    PHASE                           OBJECTIVE}

0:00    Opening \& Context               Set professional tone, explain mission\par

0:15    System Architecture Lecture     Teach power flow, design principles\par

0:45    Critical Technical Decision     Decide on Dual-BEC architecture\par

1:00    Subsystem Organization          Assign teams and ownership\par

1:15    **BREAK** (10 minutes)\par

1:25    Safety Training                 Mandatory safety protocol review\par

1:45    Hands-On Phase 1: Inventory     Check components, create tracking sheet\par

2:15    Hands-On Phase 2: Power Demo    Build power test bench, measure voltages\par

2:35    Week 1-2 Planning               Define milestones and deliverables\par

2:50    Closing \& Commitment            Confirm next meeting, sign-offs\par

\begin{itemize}
  \item SECTION 1: ROOM SETUP (Arrive 30 minutes early)
\end{itemize}

\section{WHITEBOARD LAYOUT:}

\subsection*{Title (Top): "ELECTRICAL ARCHITECTURE - Fixed-Wing UAV"}

\subsection*{Left Half: Power Flow Block Diagram}

\subsection*{(Draw during presentation - see block diagram document for reference)}

\subsection*{Battery (14.8V)}

\subsection*{↓}

\subsection*{ESC + External BEC}

\subsection*{↓}

\subsection*{Rail 1 (Flight-Critical) | Rail 2 (Servo Power)}

\subsection*{Right Half: Key Numbers}

\subsection*{- Battery: 14.8V, 2200mAh}

\subsection*{- ESC BEC: 5V @ 3A}

\subsection*{- External BEC: 5V @ 6A}

\subsection*{- Typical Draw: 1383mA}

\subsection*{- Peak Draw: 1863mA}

\subsection*{- Safety Margin: 78% (Dual-BEC)}

\subsection*{Bottom: Design Principles}

\subsection*{1. ISOLATION: Flight-critical separate from high-current loads}

\subsection*{2. REDUNDANCY: Servo BEC failure doesn't kill FC}

\subsection*{3. COMMON GROUND: All systems share ground reference}

\subsection*{4. NOISE ISOLATION: Motor wires away from sensors}

\section{TABLE LAYOUT:}

Table 1: Component Inventory \& Inspection\par

\begin{itemize}
  \item All electrical components laid out
  \item Labels for each component
  \item Checklist for verification
\end{itemize}

Table 2: Documentation Station\par

\begin{itemize}
  \item Printed documents (Power Budget, Wiring Diagram, etc.)
  \item Laptop with Google Sheet for tracking
\end{itemize}

Table 3: Power Test Bench (Later in meeting)\par

\begin{itemize}
  \item Battery (in LiPo bag)
  \item ESC
  \item Multimeter
  \item Flight Controller
  \item Servos
\end{itemize}

Table 4: Safety Equipment\par

\begin{itemize}
  \item LiPo safety bag
  \item Fire extinguisher (visible)
  \item First aid kit
  \item Safety protocol sign-off sheets
\end{itemize}

\begin{itemize}
  \item SECTION 2: OPENING (0:00 - 0:15) - Set the Tone
\end{itemize}

Stand at front of room. Make eye contact. Speak with confidence.\par

YOUR OPENING STATEMENT (Memorize or adapt):\par

"Welcome to the first electrical systems meeting for our UAV project.\par

Before we start, I want to be clear about what we're doing here. This is not a hobby\par

project. This is not something we're throwing together to fly once and forget about.\par

We are building a fixed-wing environmental monitoring UAV as a portfolio-level aerospace\par

engineering project. That means we're going to approach this like professional engineers,\par

not hobbyists.\par

\subsection{Today, we will}

\begin{itemize}
  \item Define our electrical architecture
  \item Make critical engineering decisions
  \item Assign subsystem ownership
  \item Begin hands-on validation work
\end{itemize}

\subsection{By the end of this meeting, everyone will understand}

\begin{itemize}
  \item How power flows through this aircraft
  \item What components you're responsible for
  \item What risks we need to mitigate
  \item What we're building in the next 2 weeks
\end{itemize}

We originally planned a solar-powered system, but due to budget and timeline constraints,\par

we've pivoted to a fixed-wing electric aircraft with a research payload. This is the\par

correct engineering decision, and it actually makes our electrical systems more robust.\par

Let me be direct: electrical is the critical path for this project. If we get this wrong,\par

nothing else matters. Mechanical can build the perfect airframe, software can write perfect\par

code, but if we have a power failure at altitude, we crash.\par

So we're going to do this right. We're going to use professional terminology, follow\par

systematic testing procedures, and prioritize safety above all else.\par

Questions before we dive into the architecture?"\par

\section{LEADERSHIP TONE TO ESTABLISH:}

\begin{itemize}
  \item Use these phrases:
  \item "Let me walk you through the architecture..."
  \item "From a systems engineering perspective..."
  \item "The failure mode here is..."
  \item "We need to validate this through testing..."
  \item "This gives us engineering margin because..."
  \item "Isolation prevents this failure from cascading..."
\end{itemize}

\begin{itemize}
  \item Avoid these phrases:
  \item "I think maybe we could..."
  \item "I'm not sure, but..."
  \item "Whatever you guys want to do..."
  \item "This is just a club project..."
  \item "We'll figure it out later..."
\end{itemize}

You are the expert. They need you to be confident and clear.\par

\begin{itemize}
  \item SECTION 3: SYSTEM ARCHITECTURE LECTURE (0:15 - 0:45) - Teach
\end{itemize}

OBJECTIVE: Every team member understands power flow and design principles.\par

PART 3A: Draw the Block Diagram (15 minutes)\par

Start at the top of whiteboard. Draw as you explain. Use different colors.\par

"Let's start at the source of all power: our OVONIC 4S LiPo battery."\par

[Draw battery at top]\par

"4S means 4 cells in series. Each cell is 3.7V nominal, so:\par

\begin{itemize}
  \item 3.7V $\times$ 4 = 14.8V nominal voltage
  \item Fully charged: 4.2V per cell = 16.8V total
  \item Empty: 3.3V per cell = 13.2V total (never go below this - battery damage)
  \item Capacity: 2200mAh = 2.2Ah
  \item This battery can deliver over 286 amps burst (we won't, but it can)"
\end{itemize}

[Draw line down from battery, split into two paths]\par

"Power splits into two main paths:\par

Path 1: The ESC - Electronic Speed Controller"\par

[Draw ESC on left]\par

"This does two things:\par

\begin{itemize}
  \item 1. Provides high-power 14.8V to the motor (10-30 amps)
  \item 2. Has a built-in BEC - Battery Eliminator Circuit - that gives us 5V @ 3A"
\end{itemize}

[Draw motor below ESC]\par

"Motor draws the most power - 10 to 18 amps typical, up to 30 amps full throttle.\par

That's 14.8V $\times$ 18A = 266 watts of power. This is why the motor dominates our power budget."\par

[Draw BEC output coming from ESC]\par

"The ESC's BEC gives us a 5V rail for electronics. 5V at 3A maximum.\par

This is Rail 1 - our Flight-Critical rail."\par

[Draw Rail 1 components: Flight Controller, GPS, Receiver, Arduino]\par

"Rail 1 powers:\par

\begin{itemize}
  \item MATEK F405-WMN Flight Controller: The brain. 250mA typical.
  \item BN-220 GPS: Navigation. 60mA.
  \item ELRS Receiver: Radio link to ground station. 90mA.
  \item Arduino Nano: Research payload brain. 30mA.
  \item Sensors: BME280 and MQ135. 160mA combined."
\end{itemize}

[Write currents next to each component, sum at bottom]\par

"Total: 663mA. We have 3000mA available from the BEC.\par

That leaves 2337mA margin. That's 78\% headroom. Excellent, right?"\par

[Pause for effect]\par

"Wrong. Because I haven't told you about the servos yet."\par

[Draw 4 servos]\par

"We have 4 servos:\par

\begin{itemize}
  \item 2 ailerons (roll control)
  \item 1 elevator (pitch control)
  \item 1 rudder (yaw control)
\end{itemize}

A single servo draws:\par

\begin{itemize}
  \item 10mA idle (just sitting there)
  \item 180mA moving (normal operation)
  \item 300mA peak (fast movement)
  \item 650mA stalled (control surface blocked or jammed)
\end{itemize}

In flight, all 4 servos move constantly. That's 4 $\times$ 180mA = 720mA typical, 1200mA peak."\par

[Write on whiteboard]\par

"663mA (electronics) + 720mA (servos) = 1383mA total typical draw.\par

We have 3000mA. Still okay, right? 1617mA margin."\par

[Write in red]\par

"But what if a servo stalls? Imagine a control surface gets jammed.\par

One servo tries to move, can't, draws 650mA.\par

Now: 663 + 650 + (3 $\times$ 180) = 1853mA. Still under 3A.\par

But what if TWO servos stall? Or THREE?\par

Now we're at 2500mA or more. We've exceeded the BEC capacity."\par

[Draw X through the circuit]\par

"What happens when you try to draw more current than a BEC can provide?\par

\begin{itemize}
  \item 1. Voltage sags - drops below 5V
  \item 2. Flight controller browns out - resets
  \item 3. You lose control
  \item 4. You crash"
\end{itemize}

[Let that sink in for 3 seconds]\par

"This is called a single point of failure. Everything depends on one BEC."\par

PART 3B: Introduce the Solution (10 minutes)\par

[Draw second path from battery]\par

"The solution: Add a second BEC. We call this Dual-BEC Architecture."\par

[Draw External UBEC on right side of diagram]\par

"External UBEC:\par

\begin{itemize}
  \item Input: 14.8V from battery (via XT60 splitter)
  \item Output: 5V @ 6A (twice the capacity of ESC BEC)
  \item This becomes Rail 2: Servo Power Rail"
\end{itemize}

[Draw servos connected to Rail 2 for power]\par

[Draw servo signal wires going to Flight Controller]\par

"Now the servos get power from Rail 2, signal from Flight Controller.\par

This is called power/signal separation. It's standard aerospace practice."\par

[Write on whiteboard]\par

RAIL 1 (ESC BEC, 3A):          RAIL 2 (External BEC, 6A):\par

\begin{itemize}
  \item Flight Controller              - 4$\times$ Servos
  \item GPS
  \item Receiver
  \item Arduino + Sensors
\end{itemize}

\begin{itemize}
  \item Total: 663mA                     Total: 720-1200mA
  \item Margin: 2337mA (78\%)             Margin: 4800mA (80\%)
\end{itemize}

"Now, even if ALL FOUR servos stall simultaneously (1 $\times$ 650 + 3 $\times$ 300 = 1550mA),\par

we still have margin on Rail 2, and Rail 1 is completely unaffected.\par

Flight Controller voltage: STABLE.\par

GPS lock: MAINTAINED.\par

Control link: ACTIVE.\par

If Rail 2 fails completely, the flight controller stays powered. We can't move control\par

surfaces, but we can glide to a controlled landing.\par

If Rail 1 fails, servos keep working, but we lose the brain - catastrophic.\par

So we protect Rail 1 at all costs."\par

[Circle on whiteboard]\par

"This is ISOLATION. High-current loads cannot affect flight-critical systems."\par

PART 3C: Design Principles (5 minutes)\par

[Write on whiteboard, explain each]\par

"Five design principles guide our electrical architecture:\par

\section{1. ISOLATION}

\begin{itemize}
  \item Flight-critical systems are isolated from high-current loads.
  \item Servo power rail is separate from FC power rail.
\end{itemize}

\section{2. REDUNDANCY}

\begin{itemize}
  \item If servo BEC fails, FC stays alive.
  \item If FC fails... we're in trouble, so we protect FC at all costs.
\end{itemize}

\section{3. COMMON GROUND}

\begin{itemize}
  \item All systems share the same ground reference.
  \item Battery negative is the master ground for everything.
  \item This prevents ground loops and signal noise.
\end{itemize}

\section{4. NOISE ISOLATION}

\begin{itemize}
  \item Motor wires carry high-frequency noise from ESC switching.
\end{itemize}

\subsection{Keep motor wires AWAY from}

\begin{itemize}
  \item GPS antenna
  \item Receiver antenna
  \item Analog sensors (MQ135)
  \item Twist motor wires together to reduce electromagnetic interference.
\end{itemize}

\section{5. POWER MARGIN}

\begin{itemize}
  \item We size everything with headroom.
  \item If a BEC can provide 3A, we only use 1.4A typical.
  \item This accounts for component tolerance, aging, and unexpected spikes.
  \item Real aerospace systems use 50-100\% margin. We're doing the same."
\end{itemize}

[Turn to team]\par

"Questions on the architecture before we make our technical decision?"\par

[Take questions. This is important - make sure everyone understands before moving forward]\par

\begin{itemize}
  \item SECTION 4: CRITICAL TECHNICAL DECISION (0:45 - 1:00) - Decide
\end{itemize}

OBJECTIVE: Team commits to Dual-BEC architecture (Option B).\par

"We have a decision to make. I'm going to present three options, then we'll discuss and vote."\par

[Refer to Power Budget document, Section 5]\par

OPTION A: Single BEC Configuration\par

"Use ESC BEC for everything. Simplest wiring.\par

\subsection{Pros}

\begin{itemize}
  \item Simple
  \item Minimal wiring
  \item One less component
\end{itemize}

\subsection{Cons}

\begin{itemize}
  \item Only 1137mA margin under peak load (38\%)
  \item Servo stall can exceed BEC capacity
  \item Single point of failure
  \item Risk of brownout $\rightarrow$ FC reset $\rightarrow$ crash
  \item No engineering margin
\end{itemize}

My assessment: UNSAFE for portfolio-level aerospace work.\par

This is acceptable for a toy quadcopter. Not for a serious UAV."\par

OPTION B: Dual-BEC Configuration (RECOMMENDED)\par

"Add external 5V 6A UBEC for servos. Professional power distribution.\par

\subsection{Pros}

\begin{itemize}
  \item Rail 1: 78\% margin (flight-critical systems protected)
  \item Rail 2: 80\% margin (servos can handle simultaneous stall)
  \item Isolation: Servo load cannot affect FC
  \item Redundancy: Servo BEC failure doesn't kill FC
  \item Professional aerospace practice
  \item Strong portfolio demonstration
\end{itemize}

\subsection{Cons}

\begin{itemize}
  \item Slightly more complex wiring
  \item Need to purchase external BEC (\textasciitilde{}\$10-15)
  \item Additional component to integrate
\end{itemize}

Cost: \$15 for UBEC\par

Benefit: Insurance against mid-flight power failure and loss of \$500 aircraft\par

My assessment: This is the CORRECT engineering choice.\par

This is what I would do if this were my job at Lockheed Martin or Boeing."\par

OPTION C: Hybrid Configuration\par

"Similar to Option B, but Arduino powered from Rail 2 instead of Rail 1.\par

This is acceptable if you want maximum isolation. Research payload failure cannot\par

affect flight systems. But adds complexity for marginal benefit."\par

\section{FACILITATE DISCUSSION}

"Thoughts? Concerns? Devil's advocate arguments?"\par

[Let team discuss for 5 minutes]\par

\subsection{Common questions you might get}

Q: "Do we really need the extra margin?"\par

A: "In engineering, you don't size for typical case. You size for worst case. What if\par

a control surface freezes due to icing? What if a servo fails and draws excessive current?\par

Margin protects us when unexpected things happen. And unexpected things ALWAYS happen."\par

Q: "Can we just test with one BEC first and add second later if needed?"\par

A: "No. If we test with insufficient power architecture and encounter brownouts during\par

testing, we've just made our debugging 10$\times$ harder. Are bugs due to software? Wiring?\par

Configuration? Or just insufficient power? We eliminate power as a variable by doing it\par

right from the start."\par

Q: "How much will the external BEC cost?"\par

A: "\$10-15 on Amazon. Hobbywing 5V 6A UBEC is \$12. Given our total project budget is\par

\$500, this is 2.4\% additional cost for significant risk mitigation. Worth it."\par

\section{VOTE \& COMMIT}

"Let's vote. Show of hands:\par

Option A (Single BEC): [Should be 0-1 hands]\par

Option B (Dual BEC):   [Should be majority]\par

Option C (Hybrid):     [Maybe 0-2 hands]"\par

\subsection{If vote is not unanimous for Option B, use this}

"I hear the concerns about added complexity. But fundamentally, this is a question of\par

risk tolerance. Option A has a >50\% chance of causing brownout during aggressive\par

maneuvering based on the power budget calculations. That's unacceptable.\par

As Electrical Lead, my recommendation is Option B. This is my formal engineering\par

decision. Does anyone have a SAFETY or TECHNICAL concern that would override this?"\par

[Address any technical concerns. If none:]\par

"Then we're implementing Option B: Dual-BEC architecture. I'll order the external UBEC\par

today. It'll arrive in 2-3 days, and we'll integrate it in Week 1."\par

[Write on whiteboard in large letters:]\par

\section{DECISION: DUAL-BEC ARCHITECTURE (OPTION B) $\checkmark$}

"This is now our baseline electrical architecture. All wiring diagrams, testing procedures,\par

and documentation will reflect this design."\par

\begin{itemize}
  \item SECTION 5: SUBSYSTEM ORGANIZATION (1:00 - 1:15) - Organize
\end{itemize}

OBJECTIVE: Assign every team member to a subsystem with clear ownership.\par

"Electrical has 4 subsystems. Each needs an OWNER and 1-2 assistants.\par

\subsection{The owner's job}

\begin{itemize}
  \item Become the expert on that subsystem
  \item Create the wiring harness or assembly
  \item Test and validate functionality
  \item Document issues and resolutions
  \item Present status at each meeting
\end{itemize}

This is your section of the electrical report. You own it."\par

\section{SUBSYSTEM 1: POWER TEAM}

\subsection{Scope}

\begin{itemize}
  \item Battery management (charging, storage, safety)
  \item XT60 connections and wiring
  \item ESC installation and wiring
  \item External BEC installation
  \item Power distribution board/harness (Rail 1 and Rail 2)
  \item Current sensing setup
  \item Power switch implementation
\end{itemize}

\subsection{Skills Needed}

\begin{itemize}
  \item Soldering (XT60 connectors, power distribution)
  \item Multimeter usage (measure voltages, current, continuity)
  \item Attention to safety (LiPo handling)
\end{itemize}

\subsection{Deliverables (Week 1-2)}

\begin{itemize}
  \item LiPo battery charged and tested (voltage check, balance check)
  \item XT60 Y-harness or splitter built (Battery $\rightarrow$ ESC + External BEC)
  \item External UBEC wired and tested (5V output verified)
  \item Power distribution harness (Rail 1 and Rail 2 wiring)
  \item Current sensing connected to FC
  \item All power connections labeled and heat-shrunk
  \item Power budget validation report (measure actual currents vs predictions)
\end{itemize}

Owner: \underline{\hspace{4.05cm}}  Assistant(s): \underline{\hspace{4.6499999999999995cm}}\par

\section{SUBSYSTEM 2: FLIGHT ELECTRONICS TEAM}

\subsection{Scope}

\begin{itemize}
  \item MATEK F405-WMN flight controller setup
  \item ESC signal wiring to FC
  \item Servo signal wiring to FC (PWM outputs)
  \item GPS module (BN-220) wiring and configuration
  \item ELRS receiver (EP2) binding and wiring
  \item Flight controller firmware configuration (iNav or ArduPilot)
  \item Servo calibration (endpoints, direction, center points)
\end{itemize}

\subsection{Skills Needed}

\begin{itemize}
  \item Flight controller software (iNav Configurator or Mission Planner)
  \item Firmware flashing (if needed)
  \item Servo calibration procedures
  \item Radio binding procedures (ELRS)
\end{itemize}

\subsection{Deliverables (Week 1-2)}

\begin{itemize}
  \item Flight controller powered and communicating via USB
  \item Firmware flashed (iNav or ArduPilot - team decision)
  \item GPS module wired, confirmed GPS lock (satellite count >8)
  \item ELRS receiver bound to transmitter (Jumper Smart ELRS TX)
  \item Receiver channels configured (verify stick inputs in FC software)
  \item ESC signal wire connected, motor direction verified (NO PROP, bench test)
  \item All 4 servos connected and responding to transmitter inputs
  \item Servo directions verified/reversed as needed (control surface check)
  \item Failsafe configured (motor cuts, servos hold last position or center)
\end{itemize}

Owner: \underline{\hspace{4.05cm}}  Assistant(s): \underline{\hspace{4.6499999999999995cm}}\par

\section{SUBSYSTEM 3: RESEARCH PAYLOAD TEAM}

\subsection{Scope}

\begin{itemize}
  \item Arduino Nano setup and programming
  \item BME280 sensor (temp/humidity/pressure) wiring and testing
  \item MQ135 sensor (air quality) wiring and testing
  \item MicroSD module wiring and data logging implementation
  \item Toggle switch integration (logging control)
  \item Data logging code development (CSV format)
  \item MQ135 warm-up procedure (24-48 hour burn-in before flight)
\end{itemize}

\subsection{Skills Needed}

\begin{itemize}
  \item Arduino programming (C++)
  \item I2C communication (BME280)
  \item Analog input reading (MQ135)
  \item SPI communication (MicroSD)
  \item File I/O (SD card library)
  \item Sensor calibration
\end{itemize}

\subsection{Deliverables (Week 1-2)}

\begin{itemize}
  \item Arduino Nano powered from 5V rail (Rail 1 or Rail 2)
  \item BME280 connected via I2C (A4/A5), test sketch confirms readings
  \item MQ135 connected to A0, analog readings logged to Serial Monitor
  \item MicroSD module connected via SPI (D10-D13), card detected and writable
  \item Toggle switch wired to D2 with internal pullup resistor
  \item Integrated data logging code:
  \item Reads BME280 (temp, humidity, pressure)
  \item Reads MQ135 (analog gas level)
  \item Timestamps each reading
  \item Writes CSV file to SD card when switch is ON
  \item File format: timestamp,temp,humidity,pressure,gas\_level
  \item MQ135 burn-in completed (powered continuously for 48 hours before first flight)
  \item Bench test: 5-minute logging session, CSV file retrieved and verified
\end{itemize}

Owner: \underline{\hspace{4.05cm}}  Assistant(s): \underline{\hspace{4.6499999999999995cm}}\par

\section{SUBSYSTEM 4: INTEGRATION \& TESTING TEAM}

\subsection{Scope}

\begin{itemize}
  \item Oversee integration of all 3 subsystems
  \item Continuity testing (verify all ground connections)
  \item Voltage measurements at all power rails
  \item Current draw measurements (clamp meter or inline sensor)
  \item Electrical noise testing (GPS lock during motor operation)
  \item Wire management (routing, zip ties, strain relief)
  \item Vibration isolation for sensitive components (GPS, FC, sensors)
  \item Final pre-flight electrical checks
\end{itemize}

\subsection{Skills Needed}

\begin{itemize}
  \item Multimeter expertise (voltage, current, continuity, resistance)
  \item Systems thinking (how subsystems interact)
  \item Troubleshooting methodology
  \item Attention to detail
\end{itemize}

\subsection{Deliverables (Week 1-2)}

\begin{itemize}
  \item Continuity check: All grounds connected (0$\Omega$ resistance from battery negative to
  \item all component grounds: ESC, FC, Arduino, servos, sensors)
  \item Voltage measurements logged:
  \item Battery voltage: 14.0-16.8V
  \item Rail 1 voltage: 5.0V $\pm$ 0.1V (no load and under load)
  \item Rail 2 voltage: 5.0V $\pm$ 0.1V (no load and under load)
  \item Current measurements logged:
  \item Rail 1 idle: \textasciitilde{}320-400mA $\checkmark$
  \item Rail 1 with servos: \textasciitilde{}1000-1400mA $\checkmark$
  \item Rail 2 idle: \textasciitilde{}40mA $\checkmark$
  \item Rail 2 with servos moving: \textasciitilde{}700-1200mA $\checkmark$
  \item Motor current: 10-18A typical (with clamp meter)
  \item Voltage sag test: Measure FC voltage during aggressive servo movement
  \item (must stay >4.8V, ideally >4.9V)
  \item Electrical noise test: GPS lock maintained during motor spin test (50\% throttle,
  \item NO PROP, satellite count must not drop)
  \item Wire routing inspection: No wires rubbing against control surfaces, sharp edges,
  \item or moving parts
  \item Vibration isolation: GPS and FC mounted with foam or rubber dampers
  \item Final pre-flight checklist created (see Wiring Diagram Section 6)
\end{itemize}

Owner: \underline{\hspace{4.05cm}}  Assistant(s): \underline{\hspace{4.6499999999999995cm}}\par

\section{ASSIGNMENT PROCESS}

"Let me explain what each subsystem does, then we'll assign teams.\par

[Explain each subsystem - read the Scope section above]\par

\subsection{Now, think about}

\begin{itemize}
  \item What skills do you have? (Soldering, programming, testing)
  \item What do you want to learn?
  \item Who do you work well with?
\end{itemize}

I need 1 owner per subsystem. Owners must commit to:\par

\begin{itemize}
  \item Attending all meetings
  \item Spending 4-6 hours per week outside meetings on your subsystem
  \item Being the go-to expert for your area
\end{itemize}

Who's interested in each subsystem?"\par

[Facilitate assignment. Aim for 2-3 people per subsystem depending on team size]\par

[Fill out Subsystem Assignment Sheet - see Section 9]\par

"Great. Owners, you are now accountable for your subsystem. Next meeting, I expect status\par

updates from each of you. If you hit a blocker, you escalate to me immediately. Don't wait\par

until the next meeting to tell me you're stuck."\par

\begin{itemize}
  \item *BREAK** (1:15 - 1:25) - 10 Minutes
\end{itemize}

"Let's take a 10-minute break. When we come back, we'll do mandatory safety training,\par

then hands-on work. Stretch, grab water, use the restroom."\par

\begin{itemize}
  \item SECTION 6: SAFETY TRAINING (1:25 - 1:45) - Critical
\end{itemize}

OBJECTIVE: Every team member understands hazards and commits to safety procedures.\par

[Stand next to safety equipment - LiPo bag, fire extinguisher, first aid kit]\par

"Before we touch any hardware, we're doing mandatory safety training.\par

\subsection{This project involves three main hazards}

1. LiPo batteries - fire/explosion risk\par

2. Propellers - laceration risk\par

3. Electrical - burns from soldering, short circuits\par

I'm going to cover each one, then you'll sign the safety protocol acknowledging you\par

understand the risks. If you're not willing to sign, you cannot participate in hands-on work.\par

No exceptions."\par

LIPO BATTERY SAFETY (10 minutes)\par

[Hold up LiPo battery IN safety bag]\par

"This is an OVONIC 4S 2200mAh LiPo battery. It stores 32.56 watt-hours of energy.\par

If mishandled, it can enter thermal runaway - an uncontrollable chemical reaction\par

that burns at over 1000$^\circ$C. LiPo fires cannot be extinguished with water. They burn for\par

10-30 minutes and release toxic smoke.\par

Google 'LiPo fire' after this meeting. Watch a few videos. Take this seriously.\par

[Refer to Safety Protocol document, Section 1]\par

\subsection{Key rules}

NEVER charge unattended. NEVER.\par

\begin{itemize}
  \item If you're charging and need to leave, stop charging, disconnect, store in bag.
\end{itemize}

NEVER discharge below 3.3V per cell.\par

\begin{itemize}
  \item Our battery management will cut power at 3.5V/cell (14.0V total).
\end{itemize}

NEVER overcharge above 4.2V per cell.\par

\begin{itemize}
  \item The B6 charger handles this automatically, but still monitor.
\end{itemize}

NEVER use a swollen or damaged battery.\par

\begin{itemize}
  \item If you see swelling or damaged wires, report immediately. Do not charge or use.
\end{itemize}

ALWAYS store in LiPo safety bag.\par

\begin{itemize}
  \item This bag contains fire if battery ignites. It's not perfect, but it's critical.
\end{itemize}

ALWAYS charge on non-flammable surface.\par

\begin{itemize}
  \item Metal, ceramic, concrete. NOT wood, carpet, or plastic.
\end{itemize}

[Demonstrate LiPo charger]\par

I'll now show you the proper charging procedure.\par

[Walk through B6 charger setup - see Safety Protocol]\par

1. Inspect battery (no swelling, no damage)\par

2. Place in safety bag on metal surface\par

3. Connect balance plug (5-pin JST-XH)\par

4. Connect main power (XT60)\par

5. Set charger: LiPo Balance Charge, 4S, 1.0A charge rate\par

6. Press START, verify 4S 14.8V, press START again\par

7. Monitor during entire charge (60-120 minutes)\par

8. When done, disconnect balance plug first, then main plug\par

9. Let cool 15 minutes before use\par

\subsection{If battery becomes hot (>50$^\circ$C), smells unusual, or swells during charging}

\begin{itemize}
  \item STOP IMMEDIATELY. Disconnect. Move outdoors if safe. Monitor for 1 hour.
\end{itemize}

[Point to fire extinguisher]\par

Fire extinguisher is here. ABC rated, works on LiPo fires (barely). If LiPo catches fire,\par

evacuate, call 911, let it burn out outdoors if possible.\par

Questions on LiPo safety?"\par

[Take questions]\par

PROPELLER SAFETY (5 minutes)\par

"A 10-inch propeller at full throttle has tip speed around 250 mph. Contact causes\par

severe lacerations or amputation. This is not an exaggeration.\par

\subsection{Key rules}

REMOVE propeller for ALL bench testing indoors.\par

\begin{itemize}
  \item We will not install the propeller until Week 2 outdoor testing.
\end{itemize}

NEVER test motor with propeller indoors.\par

\begin{itemize}
  \item No matter how careful you think you'll be. NO PROP INDOORS.
\end{itemize}

NEVER reach toward spinning propeller.\par

\begin{itemize}
  \item Obvious, but people make mistakes. Stay behind the plane of rotation.
\end{itemize}

\subsection{When we do outdoor propeller tests}

\begin{itemize}
  \item 10-meter safety perimeter
  \item Safety officer monitors perimeter
  \item Verbal warning before arming motor: 'PROPELLER TEST - CLEAR THE AREA'
  \item Wait 10 seconds for confirmation
  \item Operator stands behind aircraft
\end{itemize}

Questions on propeller safety?"\par

ELECTRICAL SAFETY (5 minutes)\par

[Show soldering iron]\par

"Soldering iron tip reaches 350$^\circ$C (660$^\circ$F). Contact causes instant severe burns.\par

\subsection{Rules}

\begin{itemize}
  \item Always use soldering iron stand
  \item Wear safety glasses (solder can splatter)
  \item Work in ventilated area (flux fumes)
  \item Wash hands after soldering (lead exposure)
  \item Unplug immediately after use
\end{itemize}

[Show multimeter]\par

"Multimeter safety:\par

\begin{itemize}
  \item Red probe to positive, black to negative/ground
  \item Use correct setting (voltage vs current)
  \item Don't probe live circuits above 50V without training
\end{itemize}

[Hold up XT60 connector]\par

"Short circuit prevention:\par

\begin{itemize}
  \item Our battery can deliver 286 amps in a short circuit
  \item That's enough to melt wires instantly and start a fire
  \item ALWAYS disconnect battery before wiring work
  \item Use heat shrink tubing on all solder joints
  \item Keep metal tools away from exposed battery terminals
\end{itemize}

Questions on electrical safety?"\par

\section{SAFETY PROTOCOL SIGN-OFF}

[Hand out Safety Protocol document]\par

"Everyone read Section 7: Safety Acknowledgment. Page 11.\par

If you understand and agree to follow these safety procedures, sign the acknowledgment.\par

If you're not willing to commit to these procedures, you can contribute to this project\par

through CAD design, documentation, or analysis, but you cannot participate in hands-on\par

electrical work.\par

Take 2 minutes to read and sign."\par

[Collect signed acknowledgments. File these - they're important if an incident occurs]\par

"Thank you. Safety is everyone's responsibility. If you see unsafe behavior, speak up.\par

I don't care if it's me doing something unsafe - call it out. No exceptions."\par

\begin{itemize}
  \item SECTION 7: HANDS-ON PHASE 1 - INVENTORY (1:45 - 2:15) - Inspect
\end{itemize}

OBJECTIVE: Verify all components present and functional. Create tracking system.\par

"Now we get hands-on. First task: verify we have all components and they're functional.\par

We're going to create an inventory spreadsheet. This is professional aerospace practice.\par

Every component is logged, inspected, and tracked."\par

[Open laptop, create Google Sheet, share with team]\par

\section{SPREADSHEET COLUMNS:}

| Component | Qty | Voltage | Current | Connector | Condition | Tested? | Assigned To | Notes |\par

[Lay out all components on Table 1]\par

"Each subsystem team, go through your components:\par

\subsection{Power Team}

\begin{itemize}
  \item OVONIC 4S LiPo battery
\end{itemize}

\section{- 40A ESC}

\begin{itemize}
  \item B6 Battery Charger
  \item LiPo Safety Bag
  \item XT60 connectors (need to verify if Y-harness included or needs to be built)
\end{itemize}

\subsection{Flight Electronics Team}

\begin{itemize}
  \item MATEK F405-WMN Flight Controller
  \item BN-220 GPS Module
  \item ELRS EP2 Receiver
  \item 4$\times$ RGBZONE ES08MA II Servos
  \item FLASH HOBBY D2830EVO Motor
  \item Folding Propeller (1060 or 1160)
\end{itemize}

\subsection{Research Payload Team}

\begin{itemize}
  \item Arduino Nano
  \item BME280 Sensor Module
  \item MQ135 Air Quality Sensor
  \item MicroSD Card Module
  \item SanDisk 32GB MicroSD Card
  \item Toggle Switch (SPST, 2-pin)
\end{itemize}

\subsection{Integration Team}

\begin{itemize}
  \item Multimeter (verify we have one)
  \item Soldering iron kit
  \item Wire (various gauges)
  \item Heat shrink tubing
  \item XT60 connectors (spares for power distribution)
  \item Zip ties, foam tape, etc.
\end{itemize}

\subsection{For each component}

1. Verify physical condition (no visible damage)\par

2. Check connector types (matches wiring diagram?)\par

3. Identify any missing components or connectors\par

4. Log in spreadsheet\par

You have 20 minutes. If you find issues, report to me immediately."\par

[Circulate among teams, answer questions, verify inventory]\par

\section{MISSING COMPONENTS DISCUSSION}

[After 20 minutes, reconvene]\par

"What are we missing?"\par

[Expected: External 5V UBEC, possibly XT60 Y-harness, possibly some connectors/wire]\par

"I'll order [missing items] today. They'll arrive in 2-3 days. In the meantime, we can\par

work on subsystems that don't require those parts."\par

[Update spreadsheet with missing items and expected arrival dates]\par

\begin{itemize}
  \item SECTION 8: HANDS-ON PHASE 2 - POWER DEMO (2:15 - 2:35) - Validate
\end{itemize}

OBJECTIVE: Build power test bench. Verify ESC BEC output. Demonstrate measurements.\par

"Power Team, you're up. We're going to build a simple test bench to verify the ESC BEC\par

output. Everyone else, watch and learn - you'll be doing similar tests with your subsystems."\par

[Clear Table 3 for power test bench]\par

STEP 1: Battery Voltage Check\par

[Multimeter in DC voltage mode, 20V range]\par

"First, we verify battery voltage. Red probe to positive, black to negative."\par

[Measure across XT60 terminals]\par

"Reading: \_\_\_\_\_\_\_V [Write on whiteboard]\par

Expected: 14.0-16.8V\par

Acceptable: >14.5V (battery has charge)\par

If <14.0V: Battery needs charging before proceeding\par

Our battery is at \_\_\_\_V, which is [good/needs charging]."\par

[If needs charging, demonstrate proper charging procedure per Safety Protocol]\par

STEP 2: ESC BEC Output Test (NO MOTOR CONNECTED)\par

"Now we'll verify the ESC's BEC output. Important: NO MOTOR connected yet.\par

We're only testing the BEC function."\par

[Connect battery to ESC input (XT60)]\par

"ESC is now powered. Listen for beeps - that's the ESC initializing."\par

[Take ESC's BEC output cable (red/black servo connector)]\par

"This is the BEC output. Red is +5V, black is ground."\par

[Multimeter: red probe to red wire, black probe to black wire]\par

"Measuring BEC output voltage..."\par

[Read voltage]\par

"Reading: \_\_\_\_\_\_\_V [Write on whiteboard]\par

Expected: 5.0V $\pm$ 0.1V\par

Acceptable range: 4.9-5.1V\par

If outside this range, ESC may be faulty.\par

Our BEC output: \_\_\_\_V $\checkmark$"\par

[Disconnect battery]\par

"This confirms Rail 1 will have stable 5V. Good."\par

STEP 3: Servo Movement Test (NO SIGNAL)\par

"Now let's test a servo. We'll connect power only, no signal."\par

[Connect one servo's red wire to BEC +5V (red), brown wire to BEC GND (black)]\par

[Do NOT connect servo signal wire yet]\par

"Servo is now powered. Move the servo arm by hand."\par

[Manually rotate servo arm]\par

"Feel that resistance? That's the servo motor. It should move smoothly with some resistance.\par

If it's completely stiff, servo is locked up - bad.\par

If it's completely loose, servo gears are stripped - bad.\par

This servo: [good/bad]"\par

[Test all 4 servos this way, log in spreadsheet]\par

STEP 4: Current Measurement Demo (If time permits)\par

"If we have a clamp meter or inline current sensor, I'll show you how to measure current\par

draw. This is how we'll validate our power budget later."\par

[If equipment available, demonstrate measuring current from BEC during servo movement]\par

"When we do full system integration next week, we'll measure:\par

\begin{itemize}
  \item Rail 1 current (should be \textasciitilde{}320mA idle, \textasciitilde{}660mA with Arduino active)
  \item Rail 2 current (should be \textasciitilde{}40mA idle, \textasciitilde{}720mA with servos moving)
  \item Battery current (should be 10-18A with motor running at cruise throttle)
\end{itemize}

This validates our power budget analysis."\par

[Disconnect everything, store battery in LiPo bag]\par

"Good work, Power Team. You've validated the first part of our power system."\par

\begin{itemize}
  \item SECTION 9: WEEK 1-2 PLANNING (2:35 - 2:50) - Plan
\end{itemize}

OBJECTIVE: Define milestones for next 2 weeks. Assign deliverables.\par

"Let's define what we're accomplishing in the next 2 weeks. We're in Validation Phase.\par

That means we validate each subsystem individually BEFORE integrating into the airframe."\par

[Write on whiteboard]\par

WEEK 1 MILESTONES (Next 7 Days)\par

\subsection{Power Team}

\begin{itemize}
  \item External UBEC arrives and tested (5V output verified)
  \item XT60 Y-harness or splitter built (Battery $\rightarrow$ ESC + UBEC)
  \item Power distribution harness designed (Rail 1 and Rail 2 wiring plan)
  \item All power connections soldered and heat-shrunk
  \item Power bench test: ESC + UBEC both outputting stable 5V simultaneously
\end{itemize}

\subsection{Flight Electronics Team}

\begin{itemize}
  \item Flight Controller connected to computer via USB
  \item Firmware flashed (team decides: iNav or ArduPilot - recommend iNav for fixed-wing)
  \item GPS module wired to FC, GPS lock confirmed outdoors (>8 satellites)
  \item Receiver bound to transmitter, channels configured in FC software
  \item Receiver channels tested: move TX sticks, verify response in FC config tool
\end{itemize}

\subsection{Research Payload Team}

\begin{itemize}
  \item Arduino Nano powered via USB, blink sketch loaded
  \item BME280 sensor wired and tested (I2C address found, readings valid)
  \item MQ135 sensor wired and tested (analog readings logged to Serial Monitor)
  \item MicroSD module wired, SD card detected, test file written successfully
  \item Basic data logging code written (reads sensors, prints to Serial)
\end{itemize}

\subsection{Integration Team}

\begin{itemize}
  \item Multimeter training completed (all team members can measure V, A, continuity)
  \item Continuity of grounds verified across all subsystems when assembled
  \item Testing checklist created for Week 2 full system integration
  \item Wire routing plan drafted (keep motor wires away from GPS/sensors)
\end{itemize}

WEEK 2 MILESTONES (Days 8-14)\par

\subsection{Power Team}

\begin{itemize}
  \item Full power system bench test: Battery $\rightarrow$ ESC + UBEC $\rightarrow$ Rail 1 + Rail 2
  \item Current measurements logged for Rail 1 and Rail 2 under load
  \item Voltage sag test: measure FC voltage during servo movement (must be >4.8V)
\end{itemize}

\subsection{Flight Electronics Team}

\begin{itemize}
  \item ESC signal wire connected to FC, motor direction verified (NO PROP, bench only)
  \item All 4 servos connected to FC PWM outputs, signal wires tested
  \item Servo configuration in FC: direction, endpoints, center points
  \item Full control surface check: move TX sticks, verify correct servo response
  \item Failsafe configured and tested (turn off TX, verify motor cuts, servos safe)
\end{itemize}

\subsection{Research Payload Team}

\begin{itemize}
  \item Integrated data logging code complete:
  \item Toggle switch controls logging on/off
  \item CSV file written to SD card
  \item Timestamps, temp, humidity, pressure, gas level all logged
  \item MQ135 burn-in started (powered continuously for 48 hours before first flight)
  \item 5-minute bench test: log data, retrieve SD card, verify CSV file readable
\end{itemize}

\subsection{Integration Team}

\begin{itemize}
  \item Full system integration: All subsystems powered simultaneously
  \item System-level current draw measured and logged
  \item GPS lock verified with motor running (50\% throttle, NO PROP)
  \item All wires routed and secured (zip ties, strain relief)
  \item Components vibration-isolated (GPS and FC on foam/rubber)
  \item Pre-flight electrical checklist completed
\end{itemize}

\section{DELIVERABLE FOR WEEK 2 MEETING}

\subsection{By next meeting (in 1 week)}

\subsection{Each subsystem owner brings}

\begin{itemize}
  \item 5-minute status presentation
  \item What you accomplished
  \item What you tested
  \item What worked
  \item What didn't work
  \item What you learned
  \item Blockers (if any)
\end{itemize}

\begin{itemize}
  \item Photos/videos of your work
  \item Show your wiring
  \item Show test results
  \item This documents the project for your portfolio
\end{itemize}

\begin{itemize}
  \item Updated component spreadsheet
  \item Mark your components as "Tested" when validated
\end{itemize}

\subsection{By Week 2 meeting}

\begin{itemize}
  \item FULL SYSTEM demonstration ready:
  \item All subsystems powered and functional on bench
  \item Data logging working
  \item Servos responding to transmitter
  \item GPS locked
  \item Ready for airframe integration
\end{itemize}

\section{MEETING SCHEDULE}

Next Meeting: [DATE, TIME, LOCATION]\par

Week 2 Meeting: [DATE, TIME, LOCATION]\par

\subsection{Between meetings}

\begin{itemize}
  \item Work on your subsystem 4-6 hours per week
  \item Communicate in team group chat (set up Discord/Slack/GroupMe)
  \item If you hit a blocker, message me IMMEDIATELY - don't wait until next meeting
  \item Document everything (photos, notes, measurements)
\end{itemize}

\begin{itemize}
  \item SECTION 10: CLOSING (2:50 - 3:00) - Commit
\end{itemize}

\section{YOUR CLOSING STATEMENT}

"That's our first electrical meeting. Let me summarize what we accomplished:\par

1. $\checkmark$ Defined electrical architecture: Dual-BEC power distribution\par

2. $\checkmark$ Made critical engineering decision: Option B architecture\par

3. $\checkmark$ Assigned subsystem ownership: Power, Flight, Research, Integration\par

4. $\checkmark$ Completed safety training: Everyone signed safety protocol\par

5. $\checkmark$ Validated component inventory: We know what we have and what we need\par

6. $\checkmark$ Demonstrated power test bench: Verified ESC BEC works\par

7. $\checkmark$ Defined Week 1-2 milestones: Clear deliverables for each team\par

This was a productive meeting. But the meeting is the easy part. The work happens between\par

meetings.\par

\subsection{You now own your subsystem. You are responsible for}

\begin{itemize}
  \item Learning how your components work
  \item Testing thoroughly before integration
  \item Documenting your work
  \item Delivering by Week 2
\end{itemize}

This is not a participation trophy project. This is real engineering work with real\par

deliverables and real deadlines. If you commit, I will support you 100\%. If you don't\par

have time, tell me now so we can adjust team assignments.\par

\subsection{Looking ahead}

\begin{itemize}
  \item Week 1-2: Subsystem validation (what we just planned)
  \item Week 3-4: Full system integration and bench testing
  \item Week 5-6: Airframe integration and ground tests
  \item Week 7-8: Flight testing and data collection
  \item Week 9-10: Analysis and final report
\end{itemize}

We're on track. You have everything you need to succeed.\par

Questions? Concerns? Confusion about your subsystem responsibilities?"\par

[Take questions]\par

"Alright. Let's make this the best aerospace project NLC STEM Club has ever done.\par

Take your assignment sheets. Take your printed documents. Study the wiring diagram and\par

power budget analysis.\par

I'll see you at the next meeting with progress updates.\par

Go build something great."\par

\begin{itemize}
  \item POST-MEETING ACTIONS (Your To-Do List)
\end{itemize}

\subsection{Within 24 hours}

\begin{itemize}
  \item Order missing components (External UBEC, connectors, etc.)
  \item Set up team communication channel (Discord, Slack, GroupMe)
  \item Send meeting summary to all attendees + faculty advisor
  \item Upload component inventory spreadsheet (Google Sheets, shared access)
  \item File safety protocol sign-off sheets (scan and keep copies)
\end{itemize}

\subsection{Within 48 hours}

\begin{itemize}
  \item Follow up with each subsystem owner individually:
  \item "How are you feeling about your subsystem? Any questions?"
  \item Review power budget analysis again (you'll be asked questions)
  \item Watch YouTube tutorials on:
  \item iNav or ArduPilot setup (whichever team chose)
  \item ELRS receiver binding
  \item Arduino sensor interfacing
\end{itemize}

\subsection{Weekly}

\begin{itemize}
  \item Check in with each subsystem owner via group chat
  \item Monitor progress on milestones
  \item Address blockers immediately
  \item Update project documentation
\end{itemize}

\subsection{Before next meeting}

\begin{itemize}
  \item Prepare meeting slides (status updates, agenda for Meeting \#2)
  \item Test any newly arrived components (External UBEC, etc.)
  \item Review subsystem deliverables to ensure they're on track
\end{itemize}

\section{APPENDICES}

Appendix A: Subsystem Assignment Sheet (FILL OUT DURING MEETING)\par

\section{POWER TEAM}

\begin{itemize}
  \item Owner: \underline{\hspace{4.95cm}}  Phone: \underline{\hspace{1.95cm}}  Email: \underline{\hspace{2.4cm}}
  \item Assistant 1: \underline{\hspace{4.05cm}}  Phone: \underline{\hspace{1.95cm}}  Email: \underline{\hspace{2.4cm}}
  \item Assistant 2: \underline{\hspace{4.05cm}}  Phone: \underline{\hspace{1.95cm}}  Email: \underline{\hspace{2.4cm}}
\end{itemize}

\section{FLIGHT ELECTRONICS TEAM}

\begin{itemize}
  \item Owner: \underline{\hspace{4.95cm}}  Phone: \underline{\hspace{1.95cm}}  Email: \underline{\hspace{2.4cm}}
  \item Assistant 1: \underline{\hspace{4.05cm}}  Phone: \underline{\hspace{1.95cm}}  Email: \underline{\hspace{2.4cm}}
  \item Assistant 2: \underline{\hspace{4.05cm}}  Phone: \underline{\hspace{1.95cm}}  Email: \underline{\hspace{2.4cm}}
\end{itemize}

\section{RESEARCH PAYLOAD TEAM}

\begin{itemize}
  \item Owner: \underline{\hspace{4.95cm}}  Phone: \underline{\hspace{1.95cm}}  Email: \underline{\hspace{2.4cm}}
  \item Assistant 1: \underline{\hspace{4.05cm}}  Phone: \underline{\hspace{1.95cm}}  Email: \underline{\hspace{2.4cm}}
  \item Assistant 2: \underline{\hspace{4.05cm}}  Phone: \underline{\hspace{1.95cm}}  Email: \underline{\hspace{2.4cm}}
\end{itemize}

\section{INTEGRATION \& TESTING TEAM}

\begin{itemize}
  \item Owner: \underline{\hspace{4.95cm}}  Phone: \underline{\hspace{1.95cm}}  Email: \underline{\hspace{2.4cm}}
  \item Assistant 1: \underline{\hspace{4.05cm}}  Phone: \underline{\hspace{1.95cm}}  Email: \underline{\hspace{2.4cm}}
  \item Assistant 2: \underline{\hspace{4.05cm}}  Phone: \underline{\hspace{1.95cm}}  Email: \underline{\hspace{2.4cm}}
\end{itemize}

ELECTRICAL LEAD (You)\par

\begin{itemize}
  \item Name: \underline{\hspace{5.1cm}}  Phone: \underline{\hspace{1.95cm}}  Email: \underline{\hspace{2.4cm}}
\end{itemize}

Appendix B: Meeting \#2 Preview (Week 1 Follow-Up)\par

\subsection{Meeting \#2 Agenda (1.5-2 hours)}

\begin{itemize}
  \item 1. Status Updates (40 min)
  \item Each subsystem: 5-min presentation + Q\&A
  \item Demonstrate progress (show working subsystem)
\end{itemize}

\begin{itemize}
  \item 2. Integration Planning (20 min)
  \item Review wiring diagram with actual components
  \item Plan wire routing and lengths
  \item Discuss mounting strategy for airframe
\end{itemize}

\begin{itemize}
  \item 3. Hands-On Work (40 min)
  \item Continue subsystem testing
  \item Begin subsystem-to-subsystem interface testing
  \item Troubleshoot any issues discovered
\end{itemize}

\begin{itemize}
  \item 4. Week 2 Planning (10 min)
  \item Define full system integration test plan
  \item Set date for full system bench test
\end{itemize}

\begin{itemize}
  \item 5. Closing (10 min)
  \item Confirm Week 2 deliverables
  \item Address any concerns
\end{itemize}

Appendix C: Leadership Tips (For You)\par

\subsection{Being a good technical leader}

\begin{itemize}
  \item BE CLEAR: Say exactly what needs to happen, by when, and who's responsible.
  \item BE CONFIDENT: You know more than your team. They're looking to you for guidance.
  \item BE SUPPORTIVE: When someone gets stuck, help them learn, don't just fix it.
  \item BE FIRM: If safety is compromised, shut it down immediately. No negotiations.
  \item BE ACCOUNTABLE: If the project fails, it's on you. Own it.
\end{itemize}

\begin{itemize}
  \item DON'T micromanage: Give ownership, then get out of the way.
  \item DON'T be a bottleneck: Empower team to make decisions within their subsystems.
  \item DON'T ignore problems: Address issues immediately, don't let them fester.
  \item DON'T take all the work: Your job is to LEAD, not do everything yourself.
\end{itemize}

\subsection{Remember}

\begin{itemize}
  \item This is a portfolio project for everyone, including you
  \item Your leadership will be more impressive than perfect technical work
  \item Employers want to see: systems thinking, communication, organization
  \item This project demonstrates all of that
\end{itemize}

You've got this.\par

\section{DOCUMENT INFORMATION}

Version: 1.0\par

Date: February 20, 2026\par

Author: Electrical Lead, NLC STEM Club\par

Purpose: Meeting \#1 Agenda \& Leadership Guide\par

Status: Ready for Meeting\par

\section{PRINT THIS DOCUMENT. BRING IT TO THE MEETING. FOLLOW IT.}

Good luck. You're going to do great.\par

\end{document}
